\documentclass[11pt]{article}
\usepackage[margin=1in]{geometry}
\usepackage{amsmath,amssymb,amsthm}
\usepackage[T1]{fontenc}
\newtheorem{theorem}{Theorem}
\title{Infinitely many affine-inequivalent quadratic PCF maps\\
  A disproof of the first assertion of conjecture 00000007755}
\author{earthking11}
\date{}

\begin{document}
\maketitle

\paragraph{Correction to PR \#826.}
The previous submission used $(X-a)^2+a$. Although every member is PCF,
these polynomials all lie in the affine-conjugacy class of $X^2$. That family
does not give distinct parameters in the intended moduli space. We replace it
with a construction of infinitely many PCF maps in the normalized quadratic
family $F_c(X)=X^2+c$.

Put $K=\overline{\mathbb Q}$. A polynomial over $K$ is PCF if every root of
its formal derivative in $K$ has a preperiodic forward orbit. The only
critical point of $F_c$ is $0$, because $F_c'(X)=2X$.

\begin{theorem}
There are infinitely many parameters $c\in K$ for which $F_c$ is PCF.
The corresponding maps are pairwise inequivalent under affine conjugacy.
Consequently the finiteness assertion of conjecture 00000007755 is false
already for degree $2$, even in the space of affine-conjugacy classes.
\end{theorem}

\begin{proof}
Work first in $K[T]$. Define the critical-orbit polynomials by
\[
P_0(T)=0,\qquad P_{n+1}(T)=P_n(T)^2+T.
\]
Induction gives $P_n(c)=F_c^{\circ n}(0)$. Define a second sequence by
\[
Q_0(T)=1,\qquad Q_{n+1}(T)=TQ_n(T)^2+1.
\]
Another induction shows $P_{n+1}(T)=TQ_n(T)$. For every $n\geq1$,
$Q_n(0)=1$ and the coefficient of $T$ in $Q_n$ is $Q_{n-1}(0)^2=1$.
Therefore $Q_n$ is nonconstant and has a root in the algebraically closed
field $K$. Such a root is nonzero because $Q_n(0)=1$.

Let $p$ be any prime. Choose a nonzero root $c_p$ of $Q_{p-1}$. Then
$P_p(c_p)=c_pQ_{p-1}(c_p)=0$, so $0$ is periodic under $F_{c_p}$ with
period dividing $p$. Its period is not $1$, since $F_{c_p}(0)=c_p\neq0$.
Primality forces the minimal period to equal $p$. In particular $F_{c_p}$
is PCF, and the parameters associated with different primes are distinct.
The infinitude of primes gives infinitely many such parameters.

Finally, suppose $F_c$ and $F_d$ are affinely conjugate via
$\phi(X)=aX+b$, where $a\neq0$. The relation
$\phi\circ F_c=F_d\circ\phi$ reads
\[
a(X^2+c)+b=(aX+b)^2+d.
\]
Comparing coefficients of $X^2$ yields $a=a^2$, hence $a=1$.
Comparing coefficients of $X$ then yields $2ab=0$, hence $b=0$.
The constant terms give $c=d$. Thus distinct parameters in the normalized
family represent distinct affine-conjugacy classes.
\end{proof}

\paragraph{Formalization and scope.}
The Lean file uses $K=\operatorname{AlgebraicClosure}(\mathbb Q)$ and proves
the orbit-polynomial identity, existence of a nonzero root of every $Q_n$
with $n\geq1$, exact prime periods, infinitude of normalized PCF
parameters, and affine-conjugacy rigidity. The decisive theorems are
\begin{itemize}
  \item \texttt{TLMC7755.infinite\_normalized\_pcf\_parameters};
  \item \texttt{TLMC7755.normalized\_affine\_conjugate\_iff}.
\end{itemize}
The axiom audit lists only \texttt{propext}, \texttt{Classical.choice}, and
\texttt{Quot.sound}.
We disprove the conjecture's first assertion; we make no claim about the
separate height asymptotic.

\end{document}
